\documentclass[10pt,a4paper]{amsart}
\usepackage[margin=19mm]{geometry}
\usepackage{amsmath,amssymb,mathtools}
\usepackage[scr=rsfs,cal=euler]{mathalfa}
\usepackage{xfrac}
\usepackage[unicode,hidelinks,bookmarks=false]{hyperref}
\newcommand{\ie}{i.e.\ }
\newcommand{\cf}{cf.\ }
\newcommand{\resp}{respectively\ }
\newcommand{\Subsection}{\subsection}
\providecommand{\nobreakdash}{\nobreak\hbox{-}}
\setlength{\emergencystretch}{2em}
\multlinegap0pt
\usepackage{amssymb}
\usepackage{xcolor}
\usepackage[shortlabels]{enumitem}
\usepackage{tikz}
\usepackage{algorithm}\usepackage{algpseudocode}
\usepackage{tcolorbox}
\usepackage{bm}
\newtheorem{theorem}{Theorem}
\DeclareMathOperator{\diam}{diam}
\DeclareMathOperator{\tr}{Tr}
\title{On the $l_\infty$-analog of Algebraic Connectivity}
\author{M. Rajesh Kannan, Rahul Roy}
\date{Source: arXiv:2507.22015v2 (CC BY 4.0)}
\begin{document}
\maketitle

\noindent\emph{Source and licence.} On the $l_\infty$-analog of Algebraic Connectivity. Version arXiv:2507.22015v2 (CC BY 4.0), distributed by arXiv under the Creative Commons Attribution 4.0 International licence, http://creativecommons.org/licenses/by/4.0/, as recorded in the arXiv licence metadata for this exact version and shown on the abstract page https://arxiv.org/abs/2507.22015v2. Full licence notice: \texttt{licenses/arxiv-2507.22015-CC-BY-4.0.txt}.\par\smallskip

\noindent\textit{Editorial scope.} Counted unit: the article's theorem theorem (source0.tex lines 170--228). The article's complete original proof of it is delivered with it (source0.tex lines 244--262, one \texttt{proof} environment of 19 lines, which the article places away from the statement). No mathematics is written, repaired or re-derived: the statement and the proof are the article's own lines, in the article's own notation, and no retained line is substituted or re-flowed.  No further source-local context is needed: every cross-reference of every retained line resolves inside the two delivered spans.  Nothing was omitted from any retained span, so the package carries no omission record.  No retained line contains a citation, so the wrapper carries no bibliography and no bibliography entry.  No retained line contains a figure environment, an image-inclusion command or drawing code, so no image is delivered and the package declares no asset dependency.  Editorial formatting only, in addition: an AMS \texttt{amsart} preamble that restates the article's own declarations and macros used by the retained lines, namely the environments \texttt{theorem} on separate counters, together with the standard packages those lines need.  The retained environments print in this document as Theorem 1 in order of appearance, whereas the article prints the same items with the numbering of its own full document.  The complete native source of the article as distributed in the arXiv e-print for this exact version is bundled unchanged as \texttt{sources/182389/source0.tex}.
			\begin{theorem}\label{main theorem}
				Let $G$ be a connected graph on $n$ vertices. Then
				$$
				\gamma(G)=\frac{n}{\max\limits_{u \in V(G)}  \tr(u)}=\frac{n}{\mathcal{D}_M(G)}.
				$$
				
				\begin{proof}
					Let $x \in \mathcal{F}$. Without loss of generality, let $x_u=1$ for some $u \in V(G)$. 
					Thus, for $v \in V(G)$, we have
					$$
					x_v \geq 1- \gamma_x(G) d(u,v).
					$$
					Consequently,
					$$
					0=\sum\limits_{v \in V(G)} x_v\geq \sum\limits_{v \in V(G)} (1- \gamma_x(G) d(u,v))=n- \gamma_x(G) \tr(u).
					$$
					Therefore
					$$
					\gamma_x(G) \geq \frac{n}{\tr(u)}\geq \frac{n}{\mathcal{D}_M(G)}.
					$$
					Next, let us prove the other side of the inequality holds for some vector. Let $u \in V(G)$ with $\tr(u)=\mathcal{D}_M(G)$. Define
					$$
					y_v=1 - \frac{n}{\mathcal{D}_M(G)}d(u,v).
					$$
					Then
					$$
					\sum\limits_{v\in V(G)}y_v=\sum\limits_{v\in V(G)}\left(1 - \frac{n}{\mathcal{D}_M(G)}d(u,v)\right)=n-\frac{n}{\mathcal{D}_M(G)}\tr(u)=0.
					$$
					Let $d=\diam(G)$. Choose a diametrical pair $a,b$. For every $v$,
					\begin{align*}
						d(a,v)+d(b,v) &\geq d\\
						\sum\limits_{v \in V(G)}(d(a,v)+d(b,v))&\geq nd\\
						\tr(a)+\tr(b) &\geq nd
					\end{align*}
					Also $\mathcal{D}_M(G) \geq \tr(a)$ and $\mathcal{D}_M(G) \geq \tr(b)$. Thus
					$$
					2\mathcal{D}_M(G) \geq nd.
					$$
					Therefore 
					\begin{equation}
						\mathcal{D}_M(G) \geq \frac{nd}{2}.\notag
					\end{equation}
					Since $d(u,v) \leq d$ for any $v\in V(G)$, we have,
					$$
					y_v=1 - \frac{n}{\mathcal{D}_M(G)}d(u,v)\geq 1-\frac{nd}{\mathcal{D}_M(G)}\geq 1-2=-1.
					$$
					Thus $y \in \mathcal{F}$.\\
					Finally, for an edge $vw$
					$$
					\vert d(u,v) -d(u,w) \vert \leq 1.
					$$
					\begin{align*}
						\vert y_v -y_w\vert &=\vert 1 - \frac{n}{\mathcal{D}_M(G)}d(u,v)-1+\frac{n}{\mathcal{D}_M(G)}d(u,w)\vert\\
						&\leq \frac{n}{\mathcal{D}_M(G)}\vert d(u,w) -d(u,v) \vert \\
						& \leq\frac{n}{\mathcal{D}_M(G)}.
					\end{align*}
					Therefore $ \gamma_y(G) \leq \frac{n}{\mathcal{D}_M(G)}$. Thus, $$\gamma(G)=\frac{n}{\max\limits_{u \in V(G)}  \tr(u)}=\frac{n}{\mathcal{D}_M(G)}.$$
				\end{proof}
			\end{theorem}

				\begin{proof}
					Let $x$ be an $l_\infty$-Fiedler vector. Without loss of generality, let $x_u=1$ for some $u \in V(G)$. Then
					$$
					x_v \geq 1- \gamma(G) d(u,v).
					$$
					Using Theorem \ref{main theorem} and summing over all vertices, we obtain
					$$
					0\geq n-\gamma(G) \tr(u)=n(1-\frac{\tr(u)}{\mathcal{D}_M(G)})\geq 0.
					$$
					Therefore, equality holds everywhere. Hence 
					$$
					\tr(u)=\mathcal{D}_M(G)
					$$
					and 
					$$
					x_v = 1- \gamma(G) d(u,v)
					$$
					for every vertex $v \in V(G)$.
				\end{proof}

\end{document}
